\scanpage{21}
\mspage{12}
% archive sheet number: 21
% slanted
\margin{\uline{NB} At 1--At 3 ne concernent que $\leq$. At 4 ne fait que préciser le \unsure{passage} de $\ll$ aux strates}
% lines 1--4
\[
\text{Déf}\quad F\unsure{\gtrless|\bowtie}G\Longleftrightarrow\{F,G\}\text{ majoré }(\Longleftrightarrow F\unsure{\cup|\vee}G\text{ existe})
\]
\[
\text{\uline{At 3}}\quad F\gtrless G\Longleftrightarrow\forall X\in\widetilde{\deleted{\unsure{X|F}}F},\ Y\in\widetilde{G},\ X\gtrless Y
\]
(NB $\Rightarrow$ est déjà connu)
\[
\text{\uline{At 4}}\quad F\ll G\Longleftrightarrow\forall X\in\widetilde{F},\ \exists Y\in\widetilde{G},\ X\ll Y
\]
% lines 5--7
\[
\text{\uline{At 5}}\quad\text{a)}\quad F\mid\circ\mid G\Longrightarrow F,G\in M.
\]
\margin{i.e. la relation $\mid\circ\mid$ est en fait une relation dans $M$}
% slanted
\margin{\unsure{relie} $\mid\circ\mid$ et $\gtrless$}
\[
\begin{aligned}
\text{b)}\quad X\deleted{\unsure{\mid\circ\mid}}\gtrless Y\Longleftrightarrow{}&\forall\unsure{X|X'}\in\widetilde{X}\setminus\widetilde{X}\cap\widetilde{Y},\ Y''\in\widetilde{Y}\setminus\widetilde{X}\cap\widetilde{Y},\;\unsure{\text{--}}\\
&\unsure{\text{--}}\quad X'\mid\circ\mid Y''
\end{aligned}
\]
\[
+\operatorname{Mag}1,\operatorname{Mag}2,\operatorname{Mag}3,\operatorname{Mag}4
\]

% lines 8--14
\deleted{\illegible{1}}
\uline{Scholie} Les deux pts de vue
\[
\begin{aligned}
\text{(i)}\quad &(M,\leq,\ll,\mid\circ\mid,\unsure{\mathscr{F}|\mathcal{F}}\subset\unsure{\mathscr{F}_M|\mathbb{F}_M|F_M|\mathscr{F}_{M^{\zeta}}|\mathscr{F}_{M^{\varphi}}|\mathscr{F}_{M^f}})\quad\text{avec}\quad
\underbrace{\operatorname{Mag}1-\operatorname{Mag}4}_{\substack{\text{concernant}\\\leq,\ll,\mid\circ\mid}}+\\
&\operatorname{Mag\,at}\\
\text{(ii)}\quad &(\mathscr{F},\leq,\ll,\mid\circ\mid)\quad\text{avec axiomes}\quad\operatorname{At}1-\operatorname{At}5\\
&+\operatorname{Mag}1-\operatorname{Mag}4
\end{aligned}
\]
sont équivalents, i.e. donnent des catégories-groupoïdes \unsure{équivalents|équivalentes}.

% lines 15--22
Le point de vue ``magasin'' est plus pratique, visiblement, pour l'énoncé et le \unsure{dévissage} des axiomes de base.
Le point de vue ``atlas'' me semble par contre plus commode pour la pratique du calcul et des raisonnements géom, une fois l'axiomatique acquise.
